% !TEX TS-program = pdflatex
% Lettre de motivation — Ingénieur Infrastructure — Tunisie Leasing et Factoring
\documentclass[a4paper,11pt]{article}
\usepackage[utf8]{inputenc}
\usepackage[T1]{fontenc}
\usepackage[scaled=0.94]{helvet}
\renewcommand{\familydefault}{\sfdefault}
\usepackage[left=2.5cm,right=2.5cm,top=2cm,bottom=2cm]{geometry}
\usepackage{xcolor}
\usepackage[hidelinks]{hyperref}
\definecolor{navy}{RGB}{23,54,93}
\pagestyle{empty}
\setlength{\parindent}{0pt}
\setlength{\parskip}{0pt}
\hypersetup{pdftitle={Lettre de motivation - Baha Eddine Belhaj Mouldi - Ingenieur Infrastructure},pdfauthor={Baha Eddine Belhaj Mouldi}}
\begin{document}

\begin{flushright}
{\bfseries Baha Eddine Belhaj Mouldi}\\
Tunis, Tunisie\\
+216 55 128 982\\
\href{mailto:bahaeddine.belhajmouldi@gmail.com}{bahaeddine.belhajmouldi@gmail.com}\\
\href{https://linkedin.com/in/bahamouldi}{linkedin.com/in/bahamouldi}\\
\href{https://github.com/bahamouldi}{github.com/bahamouldi}
\end{flushright}

\vspace{0.8cm}

\begin{flushleft}
Tunis, le 22 août 2026

\vspace{0.8cm}

{\bfseries TUNISIE LEASING ET FACTORING}\\
El Menzah, Tunis

\vspace{0.6cm}

{\bfseries Objet : Candidature --- Ingénieur Infrastructure (CDI)}
\end{flushleft}

\vspace{0.6cm}

Madame, Monsieur,

Votre offre d'Ingénieur Infrastructure, couvrant la conception, le déploiement et l'administration des systèmes et réseaux, correspond directement au parcours que j'ai construit ces deux dernières années, entre travaux pratiques approfondis, mission freelance et stage de fin d'études.

Pendant mon année de mission freelance chez Beehive Entreprises, j'ai installé et administré des environnements Windows Server / Active Directory (utilisateurs, groupes, GPO) et Linux, configuré des services réseau DHCP, DNS, VPN et des règles de pare-feu, et mis en œuvre des solutions de virtualisation et de sauvegarde avec Veeam pour des infrastructures clients. Cette pratique, initiée lors de mes travaux pratiques systèmes et réseaux à l'École Nationale d'Ingénieurs de Tunis, s'est prolongée durant mon stage de fin d'études, où j'ai conçu et déployé en production un cluster Kubernetes de 5 nœuds garantissant une haute disponibilité en zéro downtime --- une mise en situation concrète des enjeux de continuité de service et de reprise après incident.

La gestion des correctifs et la sécurité des accès sont également au cœur de mon expérience : dans le cadre de mon projet BeeWAF, j'ai assuré le virtual patching de 37 CVE et maintenu une conformité continue sur 7 référentiels de sécurité (OWASP, PCI DSS, GDPR, NIST, ISO 27001, CIS, SOC 2), tout en supervisant les performances via Prometheus, Grafana et la stack ELK --- une approche que je souhaite mettre au service de la fiabilité de votre infrastructure.

Rigoureux, autonome et habitué au travail en équipe --- notamment lors de la coordination d'un groupe de 8 personnes sur un projet technique --- je porte une attention constante à la communication claire des enjeux techniques, qu'il s'agisse de documenter une infrastructure ou d'échanger avec des équipes non techniques.

Je me tiens à votre entière disposition pour un entretien afin de vous exposer plus en détail ma motivation et la manière dont mon profil peut contribuer à la fiabilité de votre système d'information.

\vspace{0.8cm}

\begin{flushright}
{\bfseries Baha Eddine Belhaj Mouldi}
\end{flushright}

\end{document}
